\section{Sketch of the Proof}
\label{sec:proof_sketch_alpha}
In this section, we present the proof of Theorem~\ref{thm:alpha}.  In the remaining part of this section, we use a general confidence ellipsoid  $\cE_t  \coloneqq \left \{\theta \in \RR^d: \|\theta - \hparam_t\|_{V_t}\leq \beta_t\right \}$ to represent $\cE^{\RLS}_{t,\delta^\prime, \lambdareg}$ and $\cE^{\PVT}_{t,\delta^\prime, \lambdareg}$, since the proof works for both of them. 


First note that, when $\|\hparam_t\|_{V_t}  < \beta_t$, the bound in Theorem~\ref{thm:alpha} becomes 
\$\alpha_t\leq \lambda_d^{-\frac{1}{2}}(V_t) /\lambda_1^{-\frac{1}{2}}(V_t) = \sqrt{\lambda_1(V_t)/\lambda_d(V_t)}.\$
This bound holds trivially using the fact that $\lambda_1(V_t) \leq \|x\|^2_{V_t} \leq \lambda_d(V_t) $. This is the case when the data is insufficient and the confidence interval is too large to get a non-trivial upper bound for $\alpha_t$.

In the following, without loss of generality we assume $\|\hparam_t\|_{V_t}^2  \geq \beta_t^2$. The proof decomposes into three steps. In the first two steps, as is illustrated in $\RR^2$ in Figure~\ref{fig:ellipsoid}, we cut out a special hypersurface  $\cH_t$ within  $\cE_t$ and show that for all $\theta \in \cH_t$, the corresponding optimal action $x^\star(\theta)$ has $V_t$-norm bounded from above and below. Note that the set of optimal actions for  $\theta \in \cH_t$ coincides with that for  $\theta \in \cE_t$, we get upper and lower bounds for  $V_t$-norm of all potential actions in the ellipsoid. Next, we show that upper and lower bounds for the $V_t$-norm can be converted into upper and lower bounds for the $V_t^{-1}$-norm by solving a linear programming problem. Hence, we get an upper bound for $\alpha_t$ by calculating the ratio of the upper bound to the lower bound. We sketch the proof below and postpone the detailed proof for all lemmas in this section to Appendix~\ref{sec:proof_alpha}.


\begin{figure}[htp!]
	\centering
%	\subfigure{
		\includegraphics[width=.5\textwidth]{figures/ellipsoid_7}
	%}
	\caption{Illustration of Step 1 and 2 in $\RR^2$. Orange dashed rays: rays starting from the origin might have different numbers of intersections with $\cE_t$, indicating whether the corresponding action lies in the projection of $\cE_t$ onto $\cS_{d-1}$. Blue dashed curve:  the ellipsoid with fixed $V_t$-norm  $\left\{\theta:\|\theta\|_{V_t}=\phi_t\right\}$. The intersection of this ellipsoid and $\cE_t$ has the same projection as $\cE_t$ onto $\cS_{d-1}$. }
	\label{fig:ellipsoid}
\end{figure}


In the following, we let
$
\phi_t	 \coloneqq \sqrt{ \|\hparam_t\|_{V_t}^2 - \beta_t^2 }
$, which is well-defined since we assume $\|\hparam_t\|_{V_t}  \geq \beta_t$. In geometry,  one can show that for a ray starting from the origin and intersecting $\cE_t$ only at one point, $\phi_t$ is the ${V_t}$-norm of the intersection point. 

\paragraph{Step 1. Upper Bounding the $V_t$-norm of actions.}
Our first lemma investigates such intersection and provides an upper bound for the  $V_t$-norm for the optimal actions corresponding to any $\theta\in \cE_t$. The proof is based on investigating the condition for a ray paralleling an action $x \in \cS_{d-1}$ to intersect with $\cE_t$, which means $x$ is in the projection of $\cE_t$ onto $\cS_{d-1}$ and might become the optimal action $x^\star(\theta)$. 


\begin{lemma}\label{lem:Vt_norm_upper}
For any $\theta\in \cE_t$, we have
	$
	\|x^\star(\theta )\|_{V_t} \leq  \|\hparam_t\|_{V_t^2}/\phi_t$.
\end{lemma}


   



\paragraph{Step 2. Lower Bounding the $V_t$-norm of actions.}

In order to lower bound the $V_t$-norm, we define the hypersurface  
$
\cH_t \coloneqq \mathcal{E}_t\cap\left\{\theta:\|\theta\|_{V_t}=\phi_t\right\}
$, 
i.e., the intersection of the interior of the confidence ellipsoid  $\cE_t$ and the ellipsoid $\left\{\theta:\|\theta\|_{V_t}=\phi_t\right\}$.
$\cH_t$  consists of $\theta \in \cE_t$ whose $V_t$-norm is $\phi_t$. One can check $\cH_t$ is non-empty since $ \phi_t  \hparam_t  /\|\hparam_t\|_{V_t}  \in \cE_t$, and  the projection of $\cH_t$  onto $\cS_{d-1}$ is the same as that of $\cE_t$ by convexity. Hence,  it suffices to only consider  $\theta \in \cH_t$ as the corresponding set of optimal actions coincides. A lower bound for the $V_t$-norm is given by the following lemma.



\begin{lemma}
\label{lem:Vt_norm_lower}
	For any $\theta\in \cE_t$, we have
	\$
	  \|x^\star(\theta )\|_{V_t} \geq \frac{\phi_t}{\|\hat{\theta}_t\| +  \beta_t/\lambda_d(V_t) }.\$
\end{lemma}
The proof is directly using the fact that for any $ \theta \in \cH_t$, we have $\|\theta\| \leq \|\hparam_t\| +  \beta_t/\lambda_d(V_t)  $ and $\|\theta\|_{V_t}  =  \phi_t$. Also recall  $x^\star(\theta ) = \theta /\|\theta\|$ and hence $\|x^\star(\theta )\|_{V_t} =  \|\theta\|_{V_t} / \|\theta\|$. 










\paragraph{Step 3. Bounding the $V_t^{-1}$-norm of actions}


The following lemma determines the range of action $x$'s $V_t^{-1}$-norm based on its $V_t$-norm range. It turns out that the two ranges can be related using the spectral information of the sample covariance matrix $V_t$, which is related to the shape of the confidence ellipsoid.


\begin{lemma}
\label{lem:Vinv_norm}
Let $\{\lambda_1,\lambda_2,\dots, \lambda_{N_V}\}$ be the set of distinct eigenvalues of $V$ such that $\lambda_1>\lambda_2>\dots> \lambda_{N_V}>0$. Let $x\in \cB_d$ satisfies
	$
	0< m \leq \|x\|_{V}^2 \leq M
	$. We have
	\#
	\frac{1}{ \lambda_{k}} + \frac{1}{ \lambda_{k+1}}   - \frac{M}{\lambda_{k}\lambda_{k+1} } \leq \|x\|_{V^{-1}}^2  \leq  \frac{1}{ \lambda_{1}} + \frac{1}{ \lambda_{d}}   - \frac{m}{\lambda_{1}\lambda_{d} },
	\#
	where $k$ is such that $\lambda_{k} = \max_{i \in[N_v]} \{\lambda_i  \geq M\}$.
\end{lemma}

The proof involves expressing the $V$- and $V^{-1}$-norms as weighted sums of $V$'s eigenvalues, then solving a linear programming (LP) problem constrained by the norm ranges.




By inserting the upper and lower bounds of the $V_t$-norm from Lemmas~\ref{lem:Vt_norm_upper} and \ref{lem:Vt_norm_lower} into Lemma~\ref{lem:Vinv_norm}, we finalize the proof of Theorem~\ref{thm:alpha}.















